Un fil rectiligne de longueur $L = \qty{\v{L}}{cm}$ est parcouru par un courant de \qty{\v{I}}{A}. Il est placé dans un champ magnétique uniforme de \qty{\v{B}}{mT} qui fait un angle de \qty{\v{th}}{\degree} avec le fil.

\begin{enumerate}[label=\alph*)]
    \item Déterminez la grandeur de la force magnétique exercée sur le fil.
    \item Pour quelle orientation du fil cette force serait-elle maximale, et quelle serait sa grandeur ?
    \item On utilise maintenant un fil horizontal de même longueur et de masse \qty{\v{m}}{g}, placé perpendiculairement au même champ magnétique, lui aussi horizontal. Quel courant doit circuler dans le fil pour que la force magnétique compense exactement son poids ?
\end{enumerate}
